% Part:sets-functions-relations
% Chapter: size-of-sets
% Section: pairing-alt

\documentclass[../../../include/open-logic-section]{subfiles}

\begin{document}

\olfileid{sfr}{siz}{pai-alt}

\olsection{Nog een manier om paren te nummeren}

\begin{explain}
Je kunt de paren in $\Nat^2$ ook anders opsommen. Bij sommige manieren is het
makkelijker om vanuit een nummer het bijbehorende paar terug te vinden, dus om
de inverse te bepalen. Hier is zo'n manier. Zet de opsomming nu eens niet in
een rooster. Begin met de lijst van positieve gehele getallen en zet naast elk
getal een leeg vak. Vul die vakken één voor één als volgt met paren
$\tuple{n,m}$. Begin met de paren die op de eerste plaats~$0$ hebben, dus met
de paren $\tuple{0,m}$. Zet het eerste paar, $\tuple{0,0}$, in het eerste lege
vak. Sla een leeg vak over en zet het tweede paar, $\tuple{0,1}$, in het
volgende lege vak. Sla daarna weer een leeg vak over, enzovoort. Het begin van
de opsomming is nog niet af en ziet er nu zo uit:
\[\small
\begin{array}{@{}c c c c c c c c c c c@{}}
\mathbf 1 & \mathbf 2 & \mathbf 3 & \mathbf 4 & \mathbf 5 & \mathbf 6 & \mathbf 7 & \mathbf 8 & \mathbf 9 & \mathbf{10} & \dots \\ \\
\tuple{0,0} &  & \tuple{0,1} &  & \tuple{0,2} &  & \tuple{0,3} & & \tuple{0,4} &  & \dots \\
\end{array}
\]
Doe bij de vakken die nog leeg zijn hetzelfde met de paren $\tuple{1,m}$.
Sla ook nu steeds één leeg vak over:
\[\small
\begin{array}{@{}c c c c c c c c c c c@{}}
\mathbf 1 & \mathbf 2 & \mathbf 3 & \mathbf 4 & \mathbf 5 & \mathbf 6 & \mathbf 7 & \mathbf 8 & \mathbf 9 & \mathbf{10} & \dots \\ \\
\tuple{0,0} & \tuple{1,0} & \tuple{0,1} &  & \tuple{0,2} & \tuple{1,1} & 
\tuple{0,3} & & \tuple{0,4} &  \tuple{1,2} & \dots \\
\end{array}
\]
Vul daarna op dezelfde manier de paren $\tuple{2,m}$, $\tuple{3,m}$,
enzovoort in. De volledige opsomming begint dan zo:
\[\small
\begin{array}{@{}cc c c c c c c c c c@{}}
\mathbf 1 & \mathbf 2 & \mathbf 3 & \mathbf 4 & \mathbf 5 & \mathbf 6 & \mathbf 7 & \mathbf 8 & \mathbf 9 & \mathbf{10} & \dots \\ \\
\tuple{0,0} & \tuple{1,0} & \tuple{0,1} & \tuple{2,0}  & \tuple{0,2} & 
\tuple{1,1} & \tuple{0,3} & \tuple{3,0}  & \tuple{0,4} &  \tuple{1,2} & \dots \\
\end{array}
\]
Als we de vakken in het rooster hierboven volgens deze opsomming nummeren,
ontstaat er geen nette zigzaglijn. We krijgen deze indeling:
\[
\begin{array}{ c | c | c | c | c | c | c | c }
& \mathbf 0 & \mathbf 1 & \mathbf 2 & \mathbf 3 & \mathbf 4 & \mathbf 5 & \dots \\
\hline
\mathbf 0 & 1 & 3 & 5 & 7 & 9 & 11 & \dots \\
\hline
\mathbf 1 & 2 & 6 & 10 & 14 & 18 & \dots & \dots \\
\hline
\mathbf 2 & 4 & 12 & 20 & 28 & \dots & \dots & \dots \\
\hline
\mathbf 3 & 8 & 24 & 40 & \dots & \dots & \dots & \dots \\
\hline
\mathbf 4 & 16 & 48 & \dots & \dots & \dots & \dots & \dots \\
\hline
\mathbf 5 & 32 & \dots & \dots & \dots & \dots & \dots & \dots \\
\hline
\vdots & \vdots & \vdots & \vdots & \vdots & \vdots & \vdots & \ddots\\
\end{array}
\]
De paren in rij~$0$ staan op de oneven plaatsen van onze opsomming: het paar
$\tuple{0,m}$ staat op plaats $2m+1$. De paren in de tweede rij,
$\tuple{1,m}$, staan op plaatsen met een nummer dat tweemaal een oneven getal
is, namelijk $2 \cdot (2m+1)$. De paren in de derde rij, $\tuple{2,m}$, staan
op plaatsen met een nummer dat viermaal een oneven getal is,
$4 \cdot (2m+1)$; enzovoort. Voor de rijen vermenigvuldigen we $(2m+1)$ dus
achtereenvolgens met $1$, $2$, $4$, $8$, \dots. Dat zijn precies de machten
van~$2$: $1= 2^0$, $2 = 2^1$, $4 = 2^2$, $8 = 2^3$, \dots\@ De macht boven
de 2 is steeds het eerste getal van het paar. Bij het paar $\tuple{n,m}$ is de
factor daarom $2^n$. De algemene regel is dus $2^n \cdot (2m+1)$. Met deze
regel krijgt elk paar een \emph{positief} geheel getal: $\tuple{0,0}$ staat op
plaats~$1$. Als we bij plaats~$0$ willen beginnen, trekken we~$1$ van de
uitkomst af. Dan krijgen we:
\end{explain}

\begin{ex}
De functie $h\colon \Nat^2 \to \Nat$ met de regel
\[
h(n,m) = 2^n (2m+1) - 1
\]
is een paarfunctie voor de verzameling paren natuurlijke getallen~$\Nat^2$.
\end{ex}

\begin{explain}
In onze tweede opsomming van $\Nat^2$ krijgt het paar $\tuple{0,0}$ dus de code
$h(0,0) = 2^0(2\cdot 0+1) - 1 = 0$. Het paar $\tuple{1,2}$ krijgt de code
$2^{1} \cdot (2 \cdot 2 + 1) - 1 = 2 \cdot 5 - 1 = 9$. Het paar
$\tuple{2,6}$ krijgt de code $2^{2} \cdot (2 \cdot 6 + 1) - 1 = 51$.
\end{explain}

Soms hoeven we de paren natuurlijke getallen~$\Nat^2$ alleen een code te geven.
Niet elk natuurlijk getal hoeft daarbij als code te worden gebruikt: de
codering hoeft dus niet surjectief te zijn. Als we zo'n code terugrekenen naar
een paar, lukt dat alleen voor een deel van de getallen. De inverse is dan een
partiële functie.

\begin{ex}
De functie $j\colon \Nat^2 \to \Nat^+$ met de regel
\[
j(n,m) = 2^n3^m
\]
is !!a{injective} functie $\Nat^2 \to \Nat$.
\end{ex}
\end{document}